\documentclass[12pt,a4paper]{article}
\usepackage[utf8]{inputenc}
\usepackage[T1]{fontenc}
\usepackage[magyar]{babel}
\usepackage{amsmath}
\usepackage{amssymb}
\usepackage{geometry}
\usepackage{minted}
\geometry{margin=2.5cm}
\usepackage{listings}

\title{Pruning technikák}

\date{}

\begin{document}

\maketitle

\section*{Állapottér bejárás és pruning}

Számos kombinatorikai feladat megoldható visszalépéses kereséssel (backtracking), amely az összes lehetséges részmegoldást egy állapottér-fában járja be. Az ilyen keresések azonban gyakran exponenciális számú állapotot eredményeznek, ezért a teljes bejárás nem minden esetben végrehajtható hatékonyan. Ezen segít a \textbf{pruning} (visszametszés, vágás): célja, hogy ne vizsgáljunk meg olyan ágakat a keresési fában, amelyek biztosan nem befolyásolják a végeredményt.

Néhány példa:

Vannak számok: $3, 5, 8, 10$.

Válasszunk ki néhányat úgy, hogy az összeg ne haladja meg a 12-t.

Amikor

$3 + 5 + 8 = 16$

már tudjuk, hogy

$16 > 12$

Innen teljesen felesleges továbbmenni, ezért ezt az ágat levágjuk. Ez a pruning.

Egy másik példa:
Az $N$ királynő problémában soronként egy-egy királynőt helyezünk el. Már 3 királynőt elhelyeztünk a sakktáblán, és a negyediket próbáljuk lerakni a következő sorba. Ha egy mező egy már lerakott királynővel

\begin{itemize}
    \item ugyanabban az oszlopban van,
    \item vagy ugyanazon az átlón van,
\end{itemize}

akkor azt az elhelyezést nem folytatjuk. Ez is pruning. A teljes sakktáblát nem kell végigvizsgálni.

\subsection*{Kapcsolata a backtrackinggel}
A backtracking algoritmus általában így néz ki:
\begin{verbatim}
bejár(csúcs):
    ha kész megoldás
        kiír
    különben
        minden lehetőségre (ciklus)
            ha érvényes
                bejár(következő)
\end{verbatim}

\textbf{Pruninggal:}

\begin{verbatim}
bejár(csúcs):
    ha biztosan rossz
        vissza
    ha kész megoldás
        kiír
    különben
        minden lehetőségre (ciklus)
            bejár(következő)
\end{verbatim}

A különbség csupán egyetlen plusz ellenőrzés:

\texttt{ha biztosan rossz, akkor vissza}

Ez az egy sor rengeteg időt takaríthat meg.

\textbf{Összefoglalva:}
A hatékonyság növelésének kulcstechnikája a \textbf{pruning} (visszametszés), amely során a keresés bizonyos ágait már korai fázisban elvetjük, ha azokból biztosan nem vezethetünk érvényes vagy optimális megoldáshoz. Ez lehet például egy feltétel sérülése, a célérték túllépése, vagy annak belátása, hogy a további bővítés már nem vezethet jobb eredményre.



Az első feladatban még nem használunk pruningot, a második változatában viszont igen, így jól látható a pruning jelentősége.


\section{Feladat I. -- A király őrtornyai}

A Számok Királyságának határán $N$ őrtorony áll. A király minden toronyba szeretne őrt állítani, azonban bizonyos tornyok egyszerre nem őrizhetők, mert jelzőtüzeik megzavarnák egymást.

Egy \emph{őrbeosztás} a tornyok egy olyan részhalmaza, amelyben egyetlen tiltott toronypár sem szerepel egyszerre.

Feladatod az összes lehetséges őrbeosztás előállítása.

\subsection*{Bemenet}

Az első sor két egész számot tartalmaz: $N$ és $M$,
ahol

\begin{itemize}
    \item $N$ az őrtornyok száma,
    \item $M$ a tiltott toronypárok száma.
\end{itemize}

A következő $M$ sor mindegyike két egész számot tartalmaz: $a$ és $b$, ami azt jelenti, hogy az $a$ és $b$ sorszámú torony egyszerre nem őrizhető.

\textbf{Korlátok}

\begin{itemize}
    \item $1 \le N \le 20$
    \item $0 \le M \le 100$
\end{itemize}

\subsection*{Kimenet}

Első sor:

Az érvényes őrbeosztások száma.

Ezután minden őrbeosztást külön sorba kell kiírni.

A tornyok sorszámai minden sorban növekvő sorrendben szerepeljenek.

Az üres őrbeosztást (egy üres sort) elsőként kell kiírni. A nem üres őrbeosztásokat ezután lexikografikus sorrendben kell megjeleníteni.

\subsection*{Példa}

\subsubsection*{Bemenet}

\begin{verbatim}
4 2
1 2
2 3
\end{verbatim}

\subsubsection*{Kimenet}

\begin{verbatim}
10

1
1 3
1 3 4
1 4
2
2 4
3
3 4
4
\end{verbatim}

\subsection*{Hint 1}
Használj egy logikai mátrixot (\texttt{bool tiltott[MAX][MAX]}), amely tárolja, hogy mely tornyok nem lehetnek együtt! Az aktuális őrbeosztás tornyainak sorszámait egy \texttt{vector<int> aktualis}, a kész őrbeosztásokat egy \texttt{vector<vector<int>> megoldasok} tárolja.

\subsection*{Hint 2}
Írj egy függvényt (\texttt{lehet}), amely ellenőrzi, hogy az \texttt{idx} torony felvehető-e a jelenlegi részhalmazba! Ehhez az \texttt{aktualis} elemein kell végigmenned, és megnézned, hogy valamelyik \texttt{t} elemre igaz-e a \texttt{tiltott[idx][t]}.

\subsection*{Hint 3}
A backtrack eljárásban minden toronyról döntünk, hogy bekerül-e. Két ág van: az \texttt{idx} tornyot vagy bevesszük (ha lehet), vagy kihagyjuk. Ha \texttt{idx > N}, akkor már minden toronyról döntöttünk, tehát az aktuális részhalmaz egy kész őrbeosztás.

\subsection*{Hint 4}
Az 1. ágon: \texttt{if (lehet(idx))}, azaz megnézzük, hogy az \texttt{idx} kiválasztható-e. Ha igen, betesszük, meghívjuk a \texttt{backtrack(idx + 1)}-et, majd kivesszük. A 2. ágon az \texttt{idx} tornyot NEM választjuk ki: \texttt{backtrack(idx + 1)}.

\subsection*{Hint 5}
Ez a bejárási sorrend nem lexikografikus (például az \texttt{1 3} előbb készül el, mint az \texttt{1}), ezért a végén rendezd a megoldásokat (\texttt{sort})! Az üres őrbeosztás a rendezés után automatikusan az első lesz.



\subsection*{Megoldás}
\begin{minted}[linenos,breaklines]{cpp}
#include <iostream>
#include <vector>
#include <algorithm>

using namespace std;

int N, M;

const int MAX=21; //20 torony lehet max.
bool tiltott[MAX][MAX];

vector<int> aktualis;
vector<vector<int>> megoldasok;

/*
    Eldönti, hogy az idx torony hozzáadható-e a jelenlegi őrbeosztáshoz.
*/
bool lehet(int idx)
{
    for (int i = 0; i < (int)aktualis.size(); i++)
    {
        if (tiltott[idx][aktualis[i]])
        {
            return false;
        }
    }
    return true;
}

void backtrack(int idx)
{
    if (idx > N)
    {
        megoldasok.push_back(aktualis);
        return;
    }

    // 1. ág: bevesszük
    if (lehet(idx))
    {
        aktualis.push_back(idx);
        backtrack(idx + 1);
        aktualis.pop_back();
    }

    // 2. ág: nem vesszük be
    backtrack(idx + 1);
}

int main()
{
    ios::sync_with_stdio(false);
    cin.tie(nullptr);
    cin >> N >> M;

    for (int i = 0; i < M; i++)
    {
        int a, b;
        cin >> a >> b;

        tiltott[a][b] = true;
        tiltott[b][a] = true;
    }
    backtrack(1);
    // A bejárás sorrendje nem lexikografikus (pl. az [1, 3] előbb
    // készül el, mint az [1]), ezért a végén rendezünk.
    // Az üres vector a legkisebb, így az kerül előre.
    sort(megoldasok.begin(), megoldasok.end());
    cout << megoldasok.size() << "\n";
    for (int i = 0; i < (int)megoldasok.size(); i++)
    {
        for (int j = 0; j < (int)megoldasok[i].size(); j++)
        {
            if (j > 0)
            {
                cout << " ";
            }

            cout << megoldasok[i][j];
        }

        cout << "\n";
    }
    return 0;
}
\end{minted}


\section{Feladat II. -- A király őrtornyai}

A király átnézte az összes lehetséges őrbeosztást, és azt szeretné kiválasztani, amelyben a lehető legtöbb toronyban áll őr.

Feladatod meghatározni azt az őrbeosztást,

\begin{itemize}
    \item amelyben egyetlen tiltott toronypár sem szerepel egyszerre;
    \item amely a lehető legtöbb tornyot tartalmaz.
\end{itemize}

Ha több maximális őrbeosztás is létezik, akkor a
\textbf{lexikografikusan legkisebbet} kell kiírni.

Két őrbeosztás összehasonlításakor a kisebb első eltérő toronyszámot tartalmazó tekintendő lexikografikusan kisebbnek. Ha az egyik a másik prefixe (előtagja), akkor a rövidebb tekintendő kisebbnek.

\subsection*{Bemenet}

A standard input első sorában két egész szám található: $N$ és $M$,
ahol

\begin{itemize}
    \item $N$ az őrtornyok száma,
    \item $M$ a tiltott toronypárok száma.
\end{itemize}

A következő $M$ sor mindegyike két egész számot tartalmaz: $a$ és $b$, ami azt jelenti, hogy az $a$ és $b$ sorszámú torony egyszerre nem őrizhető.

\textbf{Korlátok}

\begin{itemize}
    \item $1 \le N \le 25$
    \item $0 \le M \le 100$
\end{itemize}

\subsection*{Kimenet}

Az első sorba írd ki a maximális őrbeosztás méretét!

A második sorba írd ki a kiválasztott tornyok sorszámait növekvő sorrendben.

\subsection*{Példa}

\subsubsection*{Bemenet}

\begin{verbatim}
6 5
1 2
1 4
2 3
3 6
5 6
\end{verbatim}

\subsubsection*{Kimenet}

\begin{verbatim}
3
1 3 5
\end{verbatim}

\subsection*{Magyarázat}

A példában három toronynál nagyobb érvényes őrbeosztás nem létezik. Háromelemű érvényes őrbeosztásból négy is van: $\{1, 3, 5\}$, $\{2, 4, 5\}$, $\{2, 4, 6\}$ és $\{3, 4, 5\}$. Ezek közül az $\{1, 3, 5\}$ a lexikografikusan legkisebb, ezért ezt kell kiírni.

\subsection*{Hint 1}

Ismét három vector kellene: \texttt{forbidden} a tiltott párokra (minden toronyhoz a vele ütköző tornyok listája), valamint kettő az aktuális és a legjobb tornyokra: \texttt{current}, \texttt{best}.

\subsection*{Hint 2}

Mielőtt kiválasztasz egy új tornyot, ellenőrizd, hogy nem ütközik-e valamelyik korábban kiválasztott toronnyal.

\subsection*{Hint 3}

Alkalmazz \textbf{pruning} technikát!

Ha az aktuális részmegoldásban eddig $x$ tornyot választottál ki, még legfeljebb $r$ torony választható, és már ismert egy $k$ méretű megoldás, akkor

\[
x+r < k
\]

esetén ezen az ágon már biztosan nem található jobb megoldás.

\subsection*{Hint 4}

Ha

\[
x+r = k,
\]

akkor is elhagyható az ág, amennyiben a belőle előállítható legjobb megoldás sem lehet lexikografikusan kisebb a már megtaláltnál. (Ez egy további vágási lehetőség; a lenti megoldás csak a Hint 3-ban leírt vágást alkalmazza, egyenlőség esetén tovább keres.)

\subsection*{Hint 5}
A lexikografikus összehasonlításhoz érdemes egy külön függvényt írni! Pl.
\begin{verbatim}
bool lexicographicallySmaller(const vector<int> &a,
                              const vector<int> &b) {
    return a < b;
}
\end{verbatim}
A \texttt{vector}-ok \texttt{<} operátora éppen lexikografikusan hasonlít.

\subsection*{Megoldás}
\begin{minted}[linenos,breaklines]{cpp}
#include <bits/stdc++.h>
using namespace std;

int N, M;

vector<vector<int>> forbidden;
vector<int> current, best;

/*
Ellenőrzi, hogy az új torony felvehető-e
a jelenlegi kiválasztottak mellé.
*/
bool ok(int x) {
    for (int y : current) {
        for (int z : forbidden[y]) {
            if (z == x) return false;
        }
    }
    return true;
}

/* Lexikografikus összehasonlítás: true, ha a < b */
bool lexicographicallySmaller(const vector<int> &a, const vector<int> &b) {
    return a < b;
}

/* Pruning: megnézzük, hogy van-e esély jobb megoldásra*/
bool canImprove(int idx, int currentSize) {
    int remaining = N - idx + 1;
    return currentSize + remaining >= (int)best.size();
}

void backtrack(int idx) {
    if (idx > N) {
        if (current.size() > best.size()) {
            best = current;
        } 
        else if (current.size() == best.size()) {
            if (lexicographicallySmaller(current, best)) {
                best = current;
            }
        }
        return;
    }

    // PRUNING: ha innen már nem lehet jobb megoldás
    if (!canImprove(idx, current.size())) {
        return;
    }

    // 1. ág: NEM választjuk
    backtrack(idx + 1);

    // 2. ág: választjuk (ha lehet)
    if (ok(idx)) {
        current.push_back(idx);
        backtrack(idx + 1);
        current.pop_back();
    }
}

int main() {
    ios::sync_with_stdio(false);
    cin.tie(nullptr);
    cin >> N >> M;
    forbidden.assign(N + 1, vector<int>());

    for (int i = 0; i < M; i++) {
        int a, b;
        cin >> a >> b;
        forbidden[a].push_back(b);
        forbidden[b].push_back(a);
    }
    backtrack(1);
    cout << best.size() << "\n";

    for (int i = 0; i < (int)best.size(); i++) {
        cout << best[i];
        if (i + 1 < (int)best.size()) cout << " ";
    }
    cout << "\n";
    return 0;
}
\end{minted}


\section{Adott összegű részhalmazok -- I.}

\subsection*{Megjegyzés}
A következő feladatban a pruning különösen fontos szerepet játszik, mivel a részhalmazok generálása során jelentősen csökkenti a bejárandó állapottér méretét.

\vspace{1cm}

Adott egy $N$ darab pozitív egész számot tartalmazó tömb, valamint egy $T$ célérték. A tömbben ismétlődő elemek is előfordulhatnak.

Feladatod, hogy hatékony visszalépéses kereséssel (backtracking) keresd meg az összes olyan részhalmazt, amelyek elemeinek összege pontosan $T$.

A megoldás során különös figyelmet kell fordítani a \textbf{pruning technikák} alkalmazására, mivel a naiv bejárás időben nem fér bele.

\textbf{Fontos szabályok:}
\begin{itemize}
    \item Minden elemet legfeljebb egyszer használhatsz fel egy részhalmazban.
    \item A részhalmazok elemei mindig \textbf{növekvő (nem csökkenő) sorrendben} szerepeljenek. (Ismétlődő értékek miatt egy részhalmazban azonos értékek is lehetnek, például \texttt{1 1 6}.)
    \item Az eredményben a részhalmazokat \textbf{lexikografikus sorrendben} kell kiírni.
    \item Azonos részhalmaz nem szerepelhet kétszer, még ha az input tartalmaz ismétlődő értékeket is.
\end{itemize}

\subsection*{Bemenet}

A standard input első sora egy $N$ egész szám.

A második sor $N$ darab pozitív egész számot tartalmaz.

A harmadik sor a célösszeget, $T$-t tartalmazza.

\textbf{Korlátok:}
\begin{itemize}
    \item $1 \le N \le 100$
    \item $1 \le A[i] \le 50$
    \item $1 \le T \le 100$
    \item Garantált, hogy a megoldások száma legfeljebb $10^5$.
\end{itemize}

\subsection*{Kimenet}

Az összes megfelelő részhalmaz kiírása soronként, elemek szóközzel elválasztva.

Ha nincs megoldás, akkor:
\begin{verbatim}
Nincs megoldas
\end{verbatim}
szöveget kell kiírni.

\subsection*{Példa}

\subsubsection*{Bemenet}
\begin{verbatim}
7
10 1 2 7 6 1 5
8
\end{verbatim}

\subsubsection*{Kimenet}
\begin{verbatim}
1 1 6
1 2 5
1 7
2 6
\end{verbatim}

\subsection*{Megjegyzés a megoldáshoz}

A feladat hatékony megoldásához backtracking szükséges, amelyet több szinten is pruning technikákkal kell optimalizálni:

\begin{itemize}
    \item Ha az aktuális részösszeg meghaladja $T$-t, az ág nem folytatható.
    \item A bemenet rendezése után, ha egy elem nagyobb, mint a fennmaradó célösszeg, a további elemek vizsgálata elhagyható.
    \item Az azonos értékű elemeket egy rekurziós szinten belül ki kell hagyni a duplikációk elkerülésére.
\end{itemize}

\subsection*{Hint 1}
A megoldáshoz szükség van 3 vector-ra!
\begin{verbatim}
vector<int> a;
vector<vector<int>> result;
vector<int> current;
\end{verbatim}
A beolvasott értékek, az eredmények, és az aktuális részhalmaz elemeinek tárolására.

\subsection*{Hint 2}
A hatékony megoldás kulcsa a rendezett tömb, a növekvő indexek mentén történő bejárás, valamint a korai visszavágás (pruning) alkalmazása minden olyan ágon, amely nem vezethet érvényes megoldáshoz.

\subsection*{Hint 3}
Ha az összeg (\texttt{sum}) egyenlő $T$-vel, akkor megvan egy megoldás; ha nagyobb $T$-nél, akkor vágunk (pruning). Akkor is vágni kell (és a ciklusból kilépni), ha az \texttt{a[j] <= T - sum} feltétel hamis, mert a rendezés miatt a további elemek még nagyobbak.

\subsection*{Hint 4}
Arra is figyelni kell, hogy ne legyen benne duplikáció! Egy rekurziós szinten belül egy adott értéket csak egyszer próbáljunk ki: jegyezzük meg az előzőleg kipróbált értéket (\texttt{prev}), és ha az aktuális elem ugyanekkora, hagyjuk ki.

\subsection*{Megoldás}
\begin{minted}[linenos,breaklines]{cpp}
#include <bits/stdc++.h>
using namespace std;

int N, T;
vector<int> a;

vector<vector<int>> result;
vector<int> current;

/*
Backtracking kombinációs bejárás
i = aktuális index
sum = jelenlegi összeg
*/
void backtrack(int i, int sum) {

    if (sum == T) {
        result.push_back(current);
        return;
    }

    if (sum > T) return;

    if (i >= N) return;

    int prev = -1;

    for (int j = i; j < N; j++) {

        // duplikáció védelem egy szinten
        bool isDuplicate = false;
        if (j > i && a[j] == prev) {
            isDuplicate = true;
        }

        if (!isDuplicate) {

            // pruning: ha túl nagy, akkor innentől minden nagyobb is
            if (a[j] <= T - sum) {

                current.push_back(a[j]);
                backtrack(j + 1, sum + a[j]);
                current.pop_back();
            } 
            else {
                // rendezett tömb miatt itt már nincs értelme továbbmenni
                break;
            }
            prev = a[j];
        }
    }
}

int main() {
    ios::sync_with_stdio(false);
    cin.tie(nullptr);

    cin >> N;
    a.resize(N);
    for (int i = 0; i < N; i++) {
        cin >> a[i];
    }
    cin >> T;
    sort(a.begin(), a.end());
    backtrack(0, 0);
    if (result.empty()) {
        cout << "Nincs megoldas\n";
        return 0;
    }
    for (auto &v : result) {
        for (int i = 0; i < (int)v.size(); i++) {
            cout << v[i];
            if (i + 1 < (int)v.size()) cout << " ";
        }
        cout << "\n";
    }
    return 0;
}
\end{minted}

\section{Adott összegű részhalmazok -- II.}

Adott egy $N$ darab pozitív egész számot tartalmazó tömb, valamint egy $T$ célérték. A tömbben előfordulhatnak ismétlődő elemek is. Feladatod, hogy megtaláld az összes olyan részhalmazt (az elemek sorrendje nem számít), amelyek elemeinek összege pontosan $T$-vel egyenlő.

\textbf{Fontos szabályok:}
\begin{itemize}
    \item Minden elemet legfeljebb egyszer használhatsz fel egy részhalmazban.
    \item A talált részhalmazokat \textbf{növekvő (nem csökkenő)} sorrendben kell tárolni (azaz minden részhalmaz elemei növekvő sorrendben legyenek).
    \item A végeredményben a részhalmazok \textbf{lexikografikusan} (szótárrendszerűen) növekvő sorrendben jelenjenek meg.
    \item A megoldásod ne tartalmazzon duplikált részhalmazokat!
\end{itemize}

\subsection*{Bemenet}

A standard input első sorában egy $N$ egész szám található, amely a tömb elemeinek számát jelöli.

A második sorban $N$ darab szóközzel elválasztott pozitív egész szám található (a tömb elemei).

A harmadik sorban egyetlen $T$ egész szám található (a célösszeg).

\textbf{Korlátok:}
\begin{itemize}
    \item $1 \le N \le 100$
    \item $1 \le A[i] \le 50$
    \item $1 \le T \le 100$
    \item Garantált, hogy a megoldások száma legfeljebb $10^5$.
\end{itemize}

\subsection*{Kimenet}

A standard kimenetre írd ki az összes olyan részhalmazt, amelyek összege $T$-vel egyenlő! Minden részhalmazt külön sorba írj, az elemeket szóközzel elválasztva!

Ha nincs egyetlen megfelelő részhalmaz sem, a kimenet egyetlen sort tartalmazzon:
\begin{verbatim}
Nincs megoldas
\end{verbatim}

\subsection*{Példa}

\subsubsection*{Példa bemenet}
\begin{verbatim}
7
10 1 2 7 6 1 5
8
\end{verbatim}

\subsubsection*{Példa kimenet}
\begin{verbatim}
1 1 6
1 2 5
1 7
2 6
\end{verbatim}

\subsection*{Magyarázat}

A bemeneti tömb: $[10, 1, 2, 7, 6, 1, 5]$, a célösszeg $T = 8$.

A lehetséges részhalmazok, amelyek összege $8$:
\begin{itemize}
    \item $[1, 1, 6]$ (az első és a második 1-es, valamint a 6-os)
    \item $[1, 2, 5]$ (az egyik 1-es, a 2-es és az 5-ös)
    \item $[1, 7]$ (az egyik 1-es és a 7-es)
    \item $[2, 6]$ (a 2-es és a 6-os)
\end{itemize}

A $[10]$ nem jöhet szóba, mert $10 > 8$. A részhalmazok lexikografikusan növekvő sorrendben jelennek meg. Nincs két egyforma részhalmaz, annak ellenére, hogy két darab 1-es is van a tömbben.

\subsection*{Hint 1}
Használd a \textbf{backtracking} (visszalépéses keresés) algoritmust a részhalmazok előállításához!
\subsection*{Hint 2}
\textbf{Rendezd} a tömböt növekvő sorrendbe!
\subsection*{Hint 3}
Alkalmazz \textbf{metszést} (pruning) a hatékonyság növelésére:
\begin{itemize}
    \item Ha az aktuális részösszeg már meghaladta $T$-t, ne menj tovább ezen az ágon!
    \item Ha az aktuális elem nagyobb, mint a maradék célösszeg, a ciklusból is kiléphetsz (rendezés miatt).
    \item Az azonos értékű elemeket egy szinten belül hagyd ki, ezzel elkerülve a duplikált részhalmazokat!
\end{itemize}
\subsection*{Hint 4}
A duplikátumok kiszűréséhez egy rekurziós szinten belül tárold az előzőleg kipróbált értéket! A részhalmazok elemeit mindig növekvő sorrendben add hozzá!

\subsection*{Hint 5}
A lexikografikus sorrend garantálható, ha a tömböt rendezve indítod, és a rekurzió során az elemeket növekvő index szerint vizsgálod.

\subsection*{Megoldás}

\begin{minted}[linenos,breaklines]{cpp}
#include <iostream>
#include <vector>
#include <algorithm>

using namespace std;

void backtrack(int start, vector<int>& candidates, int target, 
               vector<int>& current, vector<vector<int>>& result) {
    // Ha elértük a célt, elmentjük a megoldást
    if (target == 0) {
        result.push_back(current);
        return;
    }
    
    // Ha túlléptük, ezt az ágat nem vizsgáljuk tovább (METSZÉS 1)
    if (target < 0) {
        return;
    }
    
    int prev = -1;
    for (int i = start; i < (int)candidates.size(); i++) {
        // Duplikátumok kihagyása (METSZÉS 2) - continue nélkül
        if (candidates[i] != prev) {
            // Ha az aktuális elem már nagyobb, mint a maradék cél,
            // a többi elem is nagyobb (rendezés miatt) -> kilépünk (METSZÉS 3)
            if (candidates[i] > target) {
                break;
            }
            
            // Választás
            current.push_back(candidates[i]);
            
            // Rekurzív hívás: i+1, mert minden elemet csak egyszer használhatunk
            backtrack(i + 1, candidates, target - candidates[i], current, result);
            
            // Visszalépés
            current.pop_back();
            
            // Az aktuális érték beállítása az előzőnek a duplikátumok kihagyásához
            prev = candidates[i];
        }
    }
}

int main() {
    ios::sync_with_stdio(false);
    cin.tie(nullptr);

    int N;
    cin >> N;
    vector<int> candidates(N);
    for (int i = 0; i < N; i++) {
        cin >> candidates[i];
    }
   
    int T;
    cin >> T;
    // Rendezés: lehetővé teszi a metszést és a lexikografikus sorrendet
    sort(candidates.begin(), candidates.end());
    vector<int> current;
    vector<vector<int>> result;
    
    backtrack(0, candidates, T, current, result);
    
    if (result.empty()) {
        cout << "Nincs megoldas\n";
    } else {
        for (const auto& subset : result) {
            for (int i = 0; i < (int)subset.size(); i++) {
                cout << subset[i];
                if (i < (int)subset.size() - 1) {
                    cout << " ";
                }
            }
            cout << "\n";
        }
    }
    return 0;
}
\end{minted}

\section{Őrzőkapuk}

A Számok Királyságában minden lehetséges sorrend egy-egy útvonalat jelent a Királyi Kertben.

A kertbe azonban csak olyan sorozatok léphetnek be, amelyeket az Őrzőkapuk engedélyeznek. Minden kapu egy szabályt figyel: azt vizsgálja, hogy az éppen kialakuló sorrend „összhangban van-e” a királyság törvényeivel.

Ahogy haladsz előre és egyre több számot fűzöl össze, minden lépésnél újabb kapuhoz érsz. Ha a kapu úgy látja, hogy a jelenlegi részleges sorrend már biztosan nem folytatható érvényes úttá, akkor nem enged tovább — vissza kell fordulnod, és másik utat keresned.

A célod, hogy megtaláld az összes olyan teljes útvonalat, amely végig átjut az összes őrzőkapun, és eléri a kert végpontját anélkül, hogy bármelyik szabályt megsértené.

A siker kulcsa nem az, hogy minden utat kipróbálj, hanem hogy felismerd, mely utak vezetnek zsákutcába még azelőtt, hogy végigjárnád őket.

\begin{itemize}
    \item Minden számot pontosan egyszer kell felhasználni.
    \item A permutációban nem állhat két azonos érték egymás mellett.
    \item A permutációban minden szomszédos párra igaz kell legyen, hogy az első elem kisebb vagy egyenlő a másodiknál VAGY a különbségük legalább $K$.
\end{itemize}

Formálisan, egy $p_1, p_2, \dots, p_N$ permutáció akkor elfogadott, ha minden $1 \le i < N$ esetén:

\[
p_i \neq p_{i+1} \quad \text{és} \quad (p_i \le p_{i+1} \;\; \lor \;\; |p_i - p_{i+1}| \ge K)
\]

\textbf{Fontos:}
\begin{itemize}
    \item Azonos értékű elemek esetén különböző indexű elemeknek számítanak (tehát külön felhasználhatók).
    \item A kimenetben a permutációkat lexikografikusan növekvő sorrendben kell kiírni.
    \item A megoldás során azonos permutáció nem szerepelhet többször.
\end{itemize}

\subsection*{Bemenet}

Az első sorban $N$ és $K$ található.

A második sor $N$ darab egész számot tartalmaz.

\textbf{Korlátok:}
\begin{itemize}
    \item $1 \le N \le 10$
    \item $1 \le K \le 100$
    \item $1 \le A[i] \le 100$
    \item Garantált, hogy a megfelelő permutációk száma legfeljebb $10^5$.
\end{itemize}

\subsection*{Kimenet}

Írd ki az összes megfelelő permutációt soronként, szóközökkel elválasztva.

Ha nincs megoldás, akkor az alábbi szöveget:
\begin{verbatim}
Nincs megoldas
\end{verbatim}

\subsection*{Példa}

\subsubsection*{Bemenet}
\begin{verbatim}
4 3
1 3 5 2
\end{verbatim}

\subsubsection*{Kimenet}
\begin{verbatim}
1 2 3 5
1 3 5 2
1 5 2 3
2 3 5 1
2 5 1 3
3 5 1 2
5 1 2 3
\end{verbatim}

\subsection*{Magyarázat}

A rendezett bemenet: $1, 2, 3, 5$, és $K = 3$. Egy szomszédos pár akkor megengedett, ha növekvő, vagy ha a csökkenés legalább $3$. A $3$-as szám után tehát csak nagyobb szám állhat (a csökkenés legfeljebb $2$ lenne), az $5$-ös után viszont az $1$ és a $2$ is következhet. Ezért például a \texttt{2 5 3 1} nem megfelelő, mert az $5 \to 3$ csökkenés csak $2$.

\subsection*{Megjegyzés}

A feladatot visszalépéses kereséssel (backtracking) oldjuk meg.

A naiv megoldás minden permutációt előállítana, és csak a végén ellenőrizné ($O(N \cdot N!)$, $N = 10$ esetén $10! = 3\,628\,800$ permutáció), ezért érdemes a \textbf{pruning alkalmazása}: a hibás részleges permutációkat azonnal elvetjük.

\subsection*{Hint 1}
A megoldás során építsd fel a permutációt lépésről lépésre, és minden lépésnél ellenőrizd a feltételeket.

\subsection*{Hint 2}
Ha egy részleges permutáció már sérti a szomszédsági feltételt, az ágat azonnal el lehet vetni (pruning).

\subsection*{Hint 3}
A bemenet rendezése segíthet a lexikografikus sorrend biztosításában.

\subsection*{Hint 4}
Az azonos értékű elemek kezelésére figyelj: bár külön indexen szerepelnek, ugyanazt a permutációt adnák többször. Ezért egy rekurziós szinten belül egy adott értéket csak egyszer próbáljunk ki (jegyezzük meg az ezen a szinten utoljára kipróbált értéket). Így nem kell a kész permutációk között utólag keresni a duplikátumokat.

\subsection*{Megoldás}

\begin{minted}[linenos,breaklines]{cpp}
#include <bits/stdc++.h>
using namespace std;

int N, K;
vector<int> a;

vector<bool> used;
vector<int> perm;
vector<vector<int>> result;

bool ok(int last, int cur) {
    if (last == -1) return true;
    if (last == cur) return false;
    if (last <= cur) return true;
    if (abs(last - cur) >= K) return true;
    return false;
}

void backtrack() {
    if ((int)perm.size() == N) {
        result.push_back(perm);
        return;
    }

    int last = perm.empty() ? -1 : perm.back();

    int prevUsedValue = INT_MIN;
    bool hasPrev = false;

    for (int i = 0; i < N; i++) {

        if (!used[i]) {

            // duplikáció pruning (azonos értékek egy szinten)
            bool duplicateSkip = false;
            if (hasPrev && a[i] == prevUsedValue) {
                duplicateSkip = true;
            }
            if (!duplicateSkip) {

                // szomszédsági feltétel pruning
                if (ok(last, a[i])) {

                    used[i] = true;
                    perm.push_back(a[i]);

                    backtrack(); //itt a rekurzió 

                    perm.pop_back(); //"visszatörlés"
                    used[i] = false;

                    prevUsedValue = a[i];
                    hasPrev = true;
                }
            }
        }
    }
}

int main() {
    ios::sync_with_stdio(false);
    cin.tie(nullptr);
    cin >> N >> K;
    a.resize(N);
    used.assign(N, false);
    for (int i = 0; i < N; i++) {
        cin >> a[i];
    }

    sort(a.begin(), a.end());
    backtrack();
    if (result.empty()) {
        cout << "Nincs megoldas\n";
        return 0;
    }
    for (auto &v : result) {
        for (int i = 0; i < (int)v.size(); i++) {
            cout << v[i] << (i + 1 < (int)v.size() ? ' ' : '\n');
        }
    }
    return 0;
}
\end{minted}


\end{document}